\documentclass[10pt,a4paper]{amsart}
\usepackage[margin=19mm]{geometry}
\usepackage{amsmath,amssymb,mathtools}
\usepackage[scr=rsfs,cal=euler]{mathalfa}
\usepackage{xfrac}
\usepackage[unicode,hidelinks,bookmarks=false]{hyperref}
\newcommand{\ie}{i.e.\ }
\newcommand{\cf}{cf.\ }
\newcommand{\resp}{respectively\ }
\newcommand{\Subsection}{\subsection}
\providecommand{\nobreakdash}{\nobreak\hbox{-}}
\setlength{\emergencystretch}{2em}
\multlinegap0pt
\usepackage{amssymb}
\usepackage{mathtools}
\usepackage[shortlabels]{enumitem}
\usepackage{xcolor}
\newtheorem{theorem}{Theorem}
\newcommand{\E}{\mathbb{E}}
\newcommand{\KL}{\mathrm{KL}}
\newcommand{\Pcal}{\mathcal{P}}
\newcommand{\Qcal}{\mathcal{Q}}
\newcommand{\REGROW}{\operatorname{REGROW}}
\newcommand{\X}{\mathsf{X}}
\newcommand{\dinfKL}{d_{\mathrm{infKL}}}
\title{Power-one sequential tests for \texorpdfstring{$\Pcal$ against $\Pcal^c$}{P against P-complement}}
\author{Ashwin Ram, Aaditya Ramdas}
\date{Source: arXiv:2604.03218v2 (CC BY 4.0)}
\begin{document}
\maketitle

\noindent\emph{Source and licence.} Power-one sequential tests for \texorpdfstring{$\Pcal$ against $\Pcal^c$}{P against P-complement}. Version arXiv:2604.03218v2 (CC BY 4.0), distributed by arXiv under the Creative Commons Attribution 4.0 International licence, http://creativecommons.org/licenses/by/4.0/, as recorded in the arXiv licence metadata for this exact version and shown on the abstract page https://arxiv.org/abs/2604.03218v2. Full licence notice: \texttt{licenses/arxiv-2604.03218-CC-BY-4.0.txt}.\par\smallskip

\noindent\textit{Editorial scope.} Counted unit: the article's theorem regrow-localweaklsc (source0.tex lines 747--757). The article's complete original proof of it is delivered with it (source0.tex lines 759--1003, one \texttt{proof} environment of 245 lines). No mathematics is written, repaired or re-derived: the statement and the proof are the article's own lines, in the article's own notation, and no retained line is substituted or re-flowed.  No further source-local context is needed: every cross-reference of every retained line resolves inside the two delivered spans.  Nothing was omitted from any retained span, so the package carries no omission record.  No retained line contains a citation, so the wrapper carries no bibliography and no bibliography entry.  No retained line contains a figure environment, an image-inclusion command or drawing code, so no image is delivered and the package declares no asset dependency.  Editorial formatting only, in addition: an AMS \texttt{amsart} preamble that restates the article's own declarations and macros used by the retained lines, namely the environments \texttt{theorem} on separate counters, together with the standard packages those lines need.  The retained environments print in this document as Theorem 1 in order of appearance, whereas the article prints the same items with the numbering of its own full document.  The complete native source of the article as distributed in the arXiv e-print for this exact version is bundled unchanged as \texttt{sources/197793/source0.tex}.
\begin{theorem}\label{regrow-localweaklsc}
Let $\X$ be Polish. Let $\varnothing\ne \Pcal\subseteq\mathcal M_1(\X)$ and denote $\Phi_{\Pcal}(R)=\inf_{P\in\Pcal}\KL(R\mid P)$. Let $\varnothing\ne \Qcal\subseteq \mathcal M_1(\X)$ and suppose that $\Phi_{\Pcal}$ is weakly lower semicontinuous at every $Q\in\Qcal$. Then, there exists a finite-valued, nondecreasing process $E^{\star}\in\mathcal E_{\infty}$ with $E_n^{\star}\ge 3/4$ for every $n$ such that for all $Q\in\Qcal$, almost surely,
\[
\liminf_{n\to\infty}\frac1n\log E_n^\star\ge \Phi_{\Pcal}(Q)
\]
And,
\[
\lim_{n\to\infty}\frac1n \E_{Q^{\infty}}[\log E_n^{\star}]=\Phi_{\Pcal}(Q)=\dinfKL(Q;\Pcal).
\]
Consequently, if $\Qcal_{\mathrm{fin}}=\{Q\in\Qcal:\Phi_{\Pcal}(Q)<\infty\}$ is nonempty, then it follows that $\REGROW_{\infty}(\Qcal_{\mathrm{fin}}; \Pcal)=0$ and $E^{\star}$ attains this value. Now if in addition we have that $\Phi_{\Pcal}>0$ for every $Q\in\Qcal$, then it follows that $E_n^{\star}\rightarrow \infty$ almost surely under $Q^{\infty}$. 
\end{theorem}

\begin{proof}
Fix a bounded-Lipschitz metric \(d_{\rm BL}\) that metrizes weak convergence on
\(\mathcal M_1(\mathcal X)\). Let \((r_k)_{k\ge1}\) be an enumeration of
\(\mathbb Q_{>0}\). For each \(k\), define
\[
  \mathcal Q_k:=\{Q\in\mathcal Q:\Phi(Q)>r_k\}.
\]
For every \(Q\in\mathcal Q_k\), weak lower semicontinuity of \(\Phi\) at \(Q\) implies that there
exists a weakly open neighborhood \(V_{k,Q}\ni Q\) such that
\[
  \Phi(R)>r_k
  \qquad\text{for all }R\in V_{k,Q}.
\]
Choose \(\delta_{k,Q}>0\) sufficiently small that the closed \(d_{\rm BL}\)-ball
\[
  C_{k,Q}:=\{R\in\mathcal M_1(\mathcal X):d_{\rm BL}(R,Q)\le \delta_{k,Q}\}
\]
is contained in \(V_{k,Q}\), and write
\[
  U_{k,Q}:=\{R\in\mathcal M_1(\mathcal X):d_{\rm BL}(R,Q)<\delta_{k,Q}\}.
\]
Then \(U_{k,Q}\) is a weakly open neighborhood of \(Q\), \(C_{k,Q}\) is weakly closed and convex,
and
\[
  \inf_{R\in C_{k,Q}}\Phi(R)\ge r_k.
\]
The family \(\{U_{k,Q}:Q\in\mathcal Q_k\}\) is an open cover of \(\mathcal Q_k\). Since
\(\mathcal X\) is Polish, \(\mathcal M_1(\mathcal X)\) with the weak topology is separable and
metrizable, hence Lindelöf. Therefore this cover has a countable subcover. We denote it by
\[
  (U_{k,j})_{j\in J_k},
\]
where \(J_k\subseteq\mathbb N\) is finite or countably infinite, and let \(C_{k,j}\) be the corresponding
closed \(d_{\rm BL}\)-ball. Thus, for all \(k\) and \(j\in J_k\),
\[
  C_{k,j}\text{ is weakly closed and convex,}
  \qquad
  \inf_{R\in C_{k,j}}\Phi(R)\ge r_k.
\]

Let
\(
  v_t:=\frac{6}{\pi^2t^2},
   t\ge1,
\)
so that \(\sum_{t\ge1}v_t=1\). Define
\[
  S:=\sum_{k\ge1}2^{-k}e^{-r_k^2}\in(0,\infty),
\]
and, for \(j\in J_k\),
\[
  w_{k,j}:=\frac{2^{-k-j}e^{-r_k^2}}{4S}.
\]
Then
\[
  \sum_{k\ge1}\sum_{j\in J_k} w_{k,j}\le \frac14.
\]
For \(n\ge0\), define
\[
  E^\star_n
  :=
  \frac34
  +
  \sum_{k\ge1}\sum_{j\in J_k}
  w_{k,j}\sum_{t=1}^n
  v_t e^{r_kt}\mathbf 1\{\widehat Q_t\in C_{k,j}\},
\]
where \(\widehat Q_t\) is the empirical measure of \(X_1,\ldots,X_t\).

First we check that \(E^\star_n<\infty\) for every deterministic \(n\). Indeed,
\[
\begin{aligned}
  \sum_{k\ge1}\sum_{j\in J_k}
  w_{k,j}\sum_{t=1}^n v_t e^{r_kt}
  &\le
  \sum_{k\ge1}\sum_{j\in J_k}
  w_{k,j} e^{r_kn}  \\
  &\le
  \frac1{4S}
  \sum_{k\ge1}2^{-k}e^{-r_k^2+r_kn}
  \sum_{j\ge1}2^{-j}.
\end{aligned}
\]
Since
\(
  -r_k^2+r_kn\le \frac{n^2}{4},
\)
the final display is at most \(e^{n^2/4}/(4S)<\infty\).

We next verify the e-process property. Fix \(P\in\mathcal P\) and an arbitrary stopping time
\(\tau\), with the convention \(E^\star_\infty:=\lim_{n\to\infty}E^\star_n\), which exists because
\(E^\star_n\) is nondecreasing in \(n\). By Tonelli's theorem,
\[
\begin{aligned}
  \mathbb E_{P^\infty}E^\star_\tau
  &=
  \frac34
  +
  \sum_{k\ge1}\sum_{j\in J_k}
  w_{k,j}
  \sum_{t\ge1}
  v_t e^{r_kt}
  P^\infty(\tau\ge t,\widehat Q_t\in C_{k,j})  \\
  &\le
  \frac34
  +
  \sum_{k\ge1}\sum_{j\in J_k}
  w_{k,j}
  \sum_{t\ge1}
  v_t e^{r_kt}
  P^t(\widehat Q_t\in C_{k,j}).
\end{aligned}
\]
Since \(C_{k,j}\) is weakly closed and convex, Csiszár's nonasymptotic Sanov bound gives
\[
  P^t(\widehat Q_t\in C_{k,j})
  \le
  \exp\left\{
    -t\inf_{R\in C_{k,j}}\KL(R\|P)
  \right\}.
\]
For every \(R\in C_{k,j}\),
\[
  \KL(R\|P)\ge \inf_{P'\in\mathcal P}\KL(R\|P')=\Phi(R),
\]
and hence
\[
  \inf_{R\in C_{k,j}}\KL(R\|P)
  \ge
  \inf_{R\in C_{k,j}}\Phi(R)
  \ge r_k.
\]
Therefore
\[
  P^t(\widehat Q_t\in C_{k,j})\le e^{-r_kt}.
\]
Substituting this bound yields
\[
\begin{aligned}
  \mathbb E_{P^\infty}E^\star_\tau
  &\le
  \frac34
  +
  \sum_{k\ge1}\sum_{j\in J_k}
  w_{k,j}
  \sum_{t\ge1}v_t e^{r_kt}e^{-r_kt}  \\
  &=
  \frac34
  +
  \sum_{k\ge1}\sum_{j\in J_k}w_{k,j}
  \le 1.
\end{aligned}
\]
Since \(P\in\mathcal P\) and \(\tau\) were arbitrary, \(E^\star\) is an e-process for
\(\mathcal P\).

We now prove the pathwise lower bound. Fix \(Q\in\mathcal Q\), and let \(r_k<\Phi(Q)\). Then
\(Q\in\mathcal Q_k\), so there is some \(j\in J_k\) such that \(Q\in U_{k,j}\). By the strong law
for empirical measures on Polish spaces, \(\widehat Q_n\Rightarrow Q\) \(Q^\infty\)-almost surely.
Since \(U_{k,j}\) is a weak neighborhood of \(Q\), it follows that, on a \(Q^\infty\)-almost sure event,
\[
  \widehat Q_n\in U_{k,j}\subseteq C_{k,j}
  \qquad\text{eventually}.
\]
Hence, eventually,
\[
  E^\star_n
  \ge
  w_{k,j}v_n e^{r_kn}.
\]
Therefore
\[
  \liminf_{n\to\infty}\frac1n\log E^\star_n
  \ge
  r_k
  +
  \liminf_{n\to\infty}\frac1n\log w_{k,j}
  +
  \liminf_{n\to\infty}\frac1n\log v_n
  =
  r_k,
\]
because \(w_{k,j}>0\) is fixed and \(n^{-1}\log v_n\to0\). Since this holds for every rational
\(r_k<\Phi(Q)\), we obtain
\[
  \liminf_{n\to\infty}\frac1n\log E^\star_n
  \ge
  \Phi(Q)
  \qquad Q^\infty\text{-a.s.}
\]
If \(\Phi(Q)=+\infty\), the same argument applied to every \(r_k\in\mathbb Q_{>0}\) gives an infinite
lower bound.

It remains to prove the expected-log optimality statement for \(Q\in\mathcal Q_{\rm fin}\). Let
\(E=(E_n)_{n\ge0}\) be any e-process for \(\mathcal P\). For every deterministic \(n\), \(E_n\) is an
\(n\)-sample e-variable, so
\[
  \mathbb E_{P^n}E_n\le1
  \qquad\text{for every }P\in\mathcal P.
\]
By the variational inequality for relative entropy, for every \(P\in\mathcal P\),
\[
  \mathbb E_{Q^n}\log E_n
  \le
  \KL(Q^n\|P^n)+\log\mathbb E_{P^n}E_n
  \le
  n\KL(Q\|P).
\]
Taking the infimum over \(P\in\mathcal P\) gives
\[
  \frac1n\mathbb E_{Q^\infty}\log E_n
  \le
  \Phi(Q).
\]
Thus no e-process can have asymptotic expected log-growth exceeding \(\Phi(Q)\) at such a \(Q\).

Applying the pathwise lower bound to \(E^\star\), and using \(E^\star_n\ge3/4\), Fatou's lemma gives
\[
\begin{aligned}
  \Phi(Q)
  &\le
  \mathbb E_{Q^\infty}
  \left[
    \liminf_{n\to\infty}
    \frac1n\log E^\star_n
  \right]  \\
  &\le
  \liminf_{n\to\infty}
  \frac1n\mathbb E_{Q^\infty}\log E^\star_n
  \le
  \limsup_{n\to\infty}
  \frac1n\mathbb E_{Q^\infty}\log E^\star_n
  \le
  \Phi(Q).
\end{aligned}
\]
Therefore
\[
  \lim_{n\to\infty}
  \frac1n\mathbb E_{Q^\infty}\log E^\star_n
  =
  \Phi(Q),
\]
as claimed.
\end{proof}

\end{document}
